% GENERATED FILE. Edit canonical lessons, then run npm run generate:books.
% canonical-source-hash: fnv1a64:ed0930bdb90ad5b3
% canonical-lessons: TE-C281-raddu-ceyu, TE-C281-samarpincu, TE-C281-nokku, TE-C281-tappipovu, TE-C281-premincu

\chapter{To cancel, To submit, To press, To get lost, To love}
\label{ch:te-to-cancel-to-submit-to-press-to-get-lost-to-love}
\hlchaptermodality{281}

\begin{chapteropening}
\textbf{By the end of this chapter:} \emph{I can say to cancel, to submit, to press, to get lost and to love.}

Complete the last lesson of chapter 281: I can say to cancel, to submit, to press, to get lost and to love.
\end{chapteropening}

\section[\texorpdfstring{\te{రద్దు} \te{చేయు} (raddu cēyu)}{రద్దు చేయు (raddu cēyu)}]{\te{రద్దు} \te{చేయు} (raddu cēyu) — to cancel}

\label{lesson:TE-C281-raddu-ceyu}



\begin{quote}
Before the new one: say the Telugu for to admit, then the Telugu for to deposit.
\end{quote}

\subsection*{You'll want to know: \te{రద్దు} \te{చేయు}}
\textbf{\te{రద్దు} \te{చేయు}} — \emph{raddu cēyu} — \textquotedblleft{}to cancel\textquotedblright{}.

\begin{grammarlens}[title={What you've built}]
One more word for this chapter.
\end{grammarlens}

\subsection*{Your turn}
\begin{itemize}
\raggedright
  \item \emph{Say it:} \emph{raddu cēyu}
  \item \emph{Say it:} \emph{raddu cēyu}, once more
  \item \emph{Say it:} say \emph{raddu cēyu}
  \item \emph{Recall:} say the Telugu for to admit, then the Telugu for to deposit, then say \emph{raddu cēyu} again
\end{itemize}

\begin{tcolorbox}[breakable,colback=teal!4,colframe=teal!35!black,title={Before you move on}]
What does \te{రద్దు} \te{చేయు} mean? (\textquotedblleft{}To cancel\textquotedblright{}.) Say it once more.
\end{tcolorbox}
\section[\texorpdfstring{\te{సమర్పించు} (samarpiñcu)}{సమర్పించు (samarpiñcu)}]{\te{సమర్పించు} (samarpiñcu) — to submit, to present}

\label{lesson:TE-C281-samarpincu}



\begin{quote}
Before the new one: say the Telugu for to deposit, then the Telugu for to cancel.
\end{quote}

\subsection*{You'll want to know: \te{సమర్పించు}}
\textbf{\te{సమర్పించు}} — \emph{samarpiñcu} — \textquotedblleft{}to submit, to present\textquotedblright{}.

\begin{grammarlens}[title={What you've built}]
One more word for this chapter.
\end{grammarlens}

\subsection*{Your turn}
\begin{itemize}
\raggedright
  \item \emph{Say it:} \emph{samarpiñcu}
  \item \emph{Say it:} \emph{samarpiñcu}, once more
  \item \emph{Say it:} say \emph{samarpiñcu}
  \item \emph{Recall:} say the Telugu for to deposit, then the Telugu for to cancel, then say \emph{samarpiñcu} again
\end{itemize}

\begin{tcolorbox}[breakable,colback=teal!4,colframe=teal!35!black,title={Before you move on}]
What does \te{సమర్పించు} mean? (\textquotedblleft{}To submit, to present\textquotedblright{}.) Say it once more.
\end{tcolorbox}
\section[\texorpdfstring{\te{నొక్కు} (nokku)}{నొక్కు (nokku)}]{\te{నొక్కు} (nokku) — to press, to push (a button)}

\label{lesson:TE-C281-nokku}



\begin{quote}
Before the new one: say the Telugu for to cancel, then the Telugu for to submit.
\end{quote}

\subsection*{You'll want to know: \te{నొక్కు}}
\textbf{\te{నొక్కు}} — \emph{nokku} — \textquotedblleft{}to press, to push (a button)\textquotedblright{}.

\begin{grammarlens}[title={What you've built}]
One more word for this chapter.
\end{grammarlens}

\subsection*{Your turn}
\begin{itemize}
\raggedright
  \item \emph{Say it:} \emph{nokku}
  \item \emph{Say it:} \emph{nokku}, once more
  \item \emph{Say it:} say \emph{nokku}
  \item \emph{Recall:} say the Telugu for to cancel, then the Telugu for to submit, then say \emph{nokku} again
\end{itemize}

\begin{tcolorbox}[breakable,colback=teal!4,colframe=teal!35!black,title={Before you move on}]
What does \te{నొక్కు} mean? (\textquotedblleft{}To press, to push (a button)\textquotedblright{}.) Say it once more.
\end{tcolorbox}
\section[\texorpdfstring{\te{తప్పిపోవు} (tappipōvu)}{తప్పిపోవు (tappipōvu)}]{\te{తప్పిపోవు} (tappipōvu) — to get lost}

\label{lesson:TE-C281-tappipovu}



\begin{quote}
Before the new one: say the Telugu for to submit, then the Telugu for to press.
\end{quote}

\subsection*{You'll want to know: \te{తప్పిపోవు}}
\textbf{\te{తప్పిపోవు}} — \emph{tappipōvu} — \textquotedblleft{}to get lost\textquotedblright{}.

\begin{grammarlens}[title={What you've built}]
One more word for this chapter.
\end{grammarlens}

\subsection*{Your turn}
\begin{itemize}
\raggedright
  \item \emph{Say it:} \emph{tappipōvu}
  \item \emph{Say it:} \emph{tappipōvu}, once more
  \item \emph{Say it:} say \emph{tappipōvu}
  \item \emph{Recall:} say the Telugu for to submit, then the Telugu for to press, then say \emph{tappipōvu} again
\end{itemize}

\begin{tcolorbox}[breakable,colback=teal!4,colframe=teal!35!black,title={Before you move on}]
What does \te{తప్పిపోవు} mean? (\textquotedblleft{}To get lost\textquotedblright{}.) Say it once more.
\end{tcolorbox}
\section[\texorpdfstring{\te{ప్రేమించు} (prēmiñcu)}{ప్రేమించు (prēmiñcu)}]{\te{ప్రేమించు} (prēmiñcu) — to love}

\label{lesson:TE-C281-premincu}



\begin{quote}
Before the new one: say the Telugu for to press, then the Telugu for to get lost.
\end{quote}

\subsection*{You'll want to know: \te{ప్రేమించు}}
\textbf{\te{ప్రేమించు}} — \emph{prēmiñcu} — \textquotedblleft{}to love\textquotedblright{}.

\begin{grammarlens}[title={What you've built}]
One more word for this chapter.
\end{grammarlens}

\subsection*{Your turn}
\begin{itemize}
\raggedright
  \item \emph{Say it:} \emph{prēmiñcu}
  \item \emph{Say it:} \emph{prēmiñcu}, once more
  \item \emph{Say it:} say \emph{prēmiñcu}
  \item \emph{Recall:} say the Telugu for to press, then the Telugu for to get lost, then say \emph{prēmiñcu} again
\end{itemize}

\begin{tcolorbox}[breakable,colback=teal!4,colframe=teal!35!black,title={Before you move on}]
What does \te{ప్రేమించు} mean? (\textquotedblleft{}To love\textquotedblright{}.) Say it once more.
\end{tcolorbox}
